\par The following notation is used for the second adjustment.

\(m',m'',m^0\) are the three angles of triangle \(m\). The angles \(m'',m'\) are opposite the transition sides, respectively, and \(m^0\) lies between them and is marked with an asterisk in the table above.

For brevity, the numbers denote the complex positions [so, for example, \(21=x_{21}+\mathrm{i}y_{21}\), where \(\mathrm{i}=\surd-1\)].

We have
\[
\begin{aligned}
1)\quad-\frac{\delta(1-21)}{1-21}
&=\operatorname{cotang}a'\,.\,\delta a'-\operatorname{cotang}a''\,.\,\delta a''\\
&\quad+\operatorname{cotang}b'\,.\,\delta b'-\operatorname{cotang}b''\,.\,\delta b''\\
&\quad+\operatorname{cotang}c'\,.\,\delta c'-\operatorname{cotang}c''\,.\,\delta c''+\ .\ .\ .\\
&\quad+\operatorname{cotang}v'\,.\,\delta v'-\operatorname{cotang}v''\,.\,\delta v''\\
&=\mu.
\end{aligned}
\]
\clearpage
\noindent\textbf{RING OF TRIANGLES AROUND OLDENBURG.}\hfill337\par\medskip
\begingroup\small\setlength{\abovedisplayskip}{5pt}\setlength{\belowdisplayskip}{5pt}
\[
\begin{aligned}
2)\quad-\delta1={}&(1-21)\GaussOldenburgBracePair{\begin{aligned}
&\operatorname{cotang}v'\,.\,\delta v'-\operatorname{cotang}v''\,.\,\delta v''\\
&+\operatorname{cotang}a'\,.\,\delta a'-\operatorname{cotang}a''\,.\,\delta a''\\
&+\operatorname{cotang}b'\,.\,\delta b'-\operatorname{cotang}b''\,.\,\delta b''\\
&+\mathrm{i}(\delta v^0+\delta a^0+\delta b^0)
\end{aligned}}\\[2pt]
&+(1-3)\{\operatorname{cotang}c'\,.\,\delta c'-\operatorname{cotang}c''\,.\,\delta c''-\mathrm{i}\,.\,\delta c^0\}\\[2pt]
&+(1-4)\{\operatorname{cotang}d'\,.\,\delta d'-\operatorname{cotang}d''\,.\,\delta d''+\mathrm{i}\,.\,\delta d^0\}\\[2pt]
&+(1-5)\{\operatorname{cotang}e'\,.\,\delta e'-\operatorname{cotang}e''\,.\,\delta e''-\mathrm{i}\,.\,\delta e^0\}\\[2pt]
&+(1-6)\GaussOldenburgBracePair{\begin{aligned}
&\operatorname{cotang}f'\,.\,\delta f'-\operatorname{cotang}f''\,.\,\delta f''\\
&+\operatorname{cotang}g'\,.\,\delta g'-\operatorname{cotang}g''\,.\,\delta g''\\
&+\operatorname{cotang}h'\,.\,\delta h'-\operatorname{cotang}h''\,.\,\delta h''\\
&+\mathrm{i}(\delta f^0+\delta g^0+\delta h^0)
\end{aligned}}\\[2pt]
&+(1-9)\{\operatorname{cotang}i'\,.\,\delta i'-\operatorname{cotang}i''\,.\,\delta i''-\mathrm{i}\,.\,\delta i^0\}\\[2pt]
&+(1-10)\{\operatorname{cotang}k'\,.\,\delta k'-\operatorname{cotang}k''\,.\,\delta k''+\mathrm{i}\,.\,\delta k^0\}\\[2pt]
&+\ .\ .\ .\ .\ .\ .\\[2pt]
&+(1-20)\{\operatorname{cotang}u'\,.\,\delta u'-\operatorname{cotang}u''\,.\,\delta u''-\mathrm{i}\,.\,\delta u^0\}.
\end{aligned}
\]
\par\smallskip Here
\[
\begin{aligned}
\delta v'&=(20\,.\,21)-(20\,.\,1)\\
\delta v''&=(1\,.\,20)-(1\,.\,21)\\
\delta v^0&=(21\,.\,1)-(21\,.\,20)\\
\delta a'&=(1\,.\,21)-(1\,.\,2)\\
\delta a''&=(2\,.\,1)-(2\,.\,21)\\
\delta a^0&=(21\,.\,2)-(21\,.\,1)
\end{aligned}
\]
\begin{center}and so on.\end{center}
\noindent[where, as before, \((1\,.\,2)\), \((1\,.\,21)\), and so on denote the direction corrections].

The condition equation from the side ratios is therefore as follows [since \(\delta21=-0{,}581-\mathrm{i}\,.\,0{,}368\), \(\delta1=+0{,}283+\mathrm{i}\,.\,0{,}794\) and \(\delta\log\mathrm{s}_{1.21}=-0{,}000\,0156\), hence]
\[
\begin{aligned}
\frac{+0{,}0000156}{0{,}43429\ .\ .\ .}\,.\,206265&=+7{,}4091=\mu\\[4pt]
-206265\,.\,\delta21&=+119839+\mathrm{i}\,.\,75905=\nu'\\[4pt]
-206265\,.\,\delta\;1&=-58373-\mathrm{i}\,.\,163774=\nu''.
\end{aligned}
\]
\noindent[the equation is:]
\endgroup
\par\vfill\noindent\small IX.\hfill43\normalsize
\clearpage
\noindent338\hfill\textbf{POSTHUMOUS PAPERS.}\par\medskip
\begingroup
\small\setlength{\tabcolsep}{5pt}\renewcommand{\arraystretch}{1.18}
\begin{center}
\begin{tabular}{|r|l||r|l|}
\hline
[ Corr.&\multicolumn{1}{c||}{Coefficient ]}&[ Corr.&\multicolumn{1}{c|}{Coefficient ]}\\
\hline
\((1\,.\,2)\)&\(-(1-21)a'\)&\((12\,.\,10)\)&\(+(1-11)l''\)\\
\((1\,.\,21)\)&\(+(1-21)(v''+a')\)&\((12\,.\,11)\)&\(-(1-11)l''-(1-12)\mathrm{i}\)\\
\((1\,.\,20)\)&\(-(1-21)v''\)&\((12\,.\,15)\)&\(-(1-15)p'+(1-12)\mathrm{i}\)\\
\((2\,.\,3)\)&\(-(1-21)b'\)&\((12\,.\,16)\)&\(+(1-15)p'\)\\
\((2\,.\,21)\)&\(+(1-21)(a''+b')\)&\((13\,.\,14)\)&\(-(1-12)n'\)\\
\((2\,.\,1)\)&\(-(1-21)a''\)&\((13\,.\,12)\)&\(+(1-12)(m''+n')\)\\
\((3\,.\,5)\)&\(-(1-4)d'\)&\((13\,.\,11)\)&\(-(1-12)m''\)\\
\((3\,.\,4)\)&\(+(1-4)d'+(1-3)\mathrm{i}\)&\((14\,.\,15)\)&\(-(1-12)o'\)\\
\((3\,.\,21)\)&\(+(1-21)b''-(1-3)\mathrm{i}\)&\((14\,.\,12)\)&\(+(1-12)(n''+o')\)\\
\((3\,.\,2)\)&\(-(1-21)b''\)&\((14\,.\,13)\)&\(-(1-12)n''\)\\
\((4\,.\,21)\)&\(+(1-3)c''\)&\((15\,.\,17)\)&\(-(1-16)q'\)\\
\((4\,.\,3)\)&\(-(1-3)c''-(1-4)\mathrm{i}\)&\((15\,.\,16)\)&\(+(1-16)q'+(1-15)\mathrm{i}\)\\
\((4\,.\,5)\)&\(-(1-5)e'+(1-4)\mathrm{i}\)&\((15\,.\,12)\)&\(+(1-12)o''-(1-15)\mathrm{i}\)\\
\((4\,.\,6)\)&\(+(1-5)e'\)&\((15\,.\,14)\)&\(-(1-12)o''\)\\
\((5\,.\,7)\)&\(-(1-6)f'\)&\((16\,.\,12)\)&\(+(1-15)p''\)\\
\((5\,.\,6)\)&\(+(1-6)f'+(1-5)\mathrm{i}\)&\((16\,.\,15)\)&\(-(1-15)p''-(1-16)\mathrm{i}\)\\
\((5\,.\,4)\)&\(+(1-4)d''-(1-5)\mathrm{i}\)&\((16\,.\,18)\)&\(-(1-18)s'+(1-16)\mathrm{i}\)\\
\((5\,.\,3)\)&\(-(1-4)d''\)&\((16\,.\,19)\)&\(+(1-18)s'\)\\
\((6\,.\,4)\)&\(+(1-5)e''\)&\((17\,.\,18)\)&\(-(1-16)r'\)\\
\((6\,.\,5)\)&\(-(1-5)e''-(1-6)\mathrm{i}\)&\((17\,.\,16)\)&\(+(1-16)(q''+r')\)\\
\((6\,.\,9)\)&\(-(1-9)i'+(1-6)\mathrm{i}\)&\((17\,.\,15)\)&\(-(1-16)q''\)\\
\((6\,.\,10)\)&\(+(1-9)i'\)&\((18\,.\,20)\)&\(-(1-19)t'\)\\
\((7\,.\,8)\)&\(-(1-6)g'\)&\((18\,.\,19)\)&\(+(1-19)t'+(1-18)\mathrm{i}\)\\
\((7\,.\,6)\)&\(+(1-6)(f''+g')\)&\((18\,.\,16)\)&\(+(1-16)r''-(1-18)\mathrm{i}\)\\
\((7\,.\,5)\)&\(-(1-6)f''\)&\((18\,.\,17)\)&\(-(1-16)r''\)\\
\((8\,.\,9)\)&\(-(1-6)h'\)&\((19\,.\,16)\)&\(+(1-18)s''\)\\
\((8\,.\,6)\)&\(+(1-6)(g''+h')\)&\((19\,.\,18)\)&\(-(1-18)s''-(1-19)\mathrm{i}\)\\
\((8\,.\,7)\)&\(-(1-6)g''\)&\((19\,.\,20)\)&\(-(1-20)u'+(1-19)\mathrm{i}\)\\
\((9\,.\,11)\)&\(-(1-10)k'\)&\((19\,.\,21)\)&\(+(1-20)u'\)\\
\((9\,.\,10)\)&\(+(1-10)k'+(1-9)\mathrm{i}\)&\((20\,.\,1)\)&\(-(1-21)v'\)\\
\((9\,.\,6)\)&\(+(1-6)h''-(1-9)\mathrm{i}\)&\((20\,.\,21)\)&\(+(1-21)v'+(1-20)\mathrm{i}\)\\
\((9\,.\,8)\)&\(-(1-6)h''\)&\((20\,.\,19)\)&\(+(1-19)t''-(1-20)\mathrm{i}\)\\
\((10\,.\,6)\)&\(+(1-9)i''\)&\((20\,.\,18)\)&\(-(1-19)t''\)\\
\((10\,.\,9)\)&\(-(1-9)i''-(1-10)\mathrm{i}\)&\((21\,.\,19)\)&\(+(1-20)u''\)\\
\((10\,.\,11)\)&\(-(1-11)l'+(1-10)\mathrm{i}\)&\((21\,.\,20)\)&\(-(1-20)u''-(1-21)\mathrm{i}\)\\
\((10\,.\,12)\)&\(+(1-11)l'\)&\((21\,.\,3)\)&\(-(1-3)c'+(1-21)\mathrm{i}\)\\
\((11\,.\,13)\)&\(-(1-12)m'\)&\((21\,.\,4)\)&\(+(1-3)c'\)\\
\((11\,.\,12)\)&\(+(1-12)m'+(1-11)\mathrm{i}\)&&\\
\((11\,.\,10)\)&\(+(1-10)k''-(1-11)\mathrm{i}\)&&\\
\((11\,.\,9)\)&\(-(1-10)k''\)&&\\
\end{tabular}
\end{center}
\endgroup
\clearpage
\noindent\textbf{RING OF TRIANGLES AROUND OLDENBURG.}\hfill339\par\medskip
\par This abbreviated table applies first to the condition equation for \(\nu''\). Namely, \(\nu''\) equals the sum of 77 terms, of which the first is
\[
=-(1-21)\operatorname{cotang}a'\,.\,(1\,.\,2),
\]
\noindent the second
\[
=+(1-21)(\operatorname{cotang}v''+\operatorname{cotang}a')(1\,.\,21)
\]
\begin{center}and so on.\end{center}

The condition equation for \(\nu'\) is obtained by replacing 1 with 21 throughout, so that terms 1, 2, 3, 4, 5, 6, 10 and 70 drop out.

Finally, the condition equation for \(\mu\) is obtained by dropping the respective first factors (\(1-21\), \(1-3\), \(1-4\), and so on) and also dropping the imaginary parts.

The following page gives the 3 condition equations (for \(\mu\) and the two parts of \(\nu'\)) in numerical form, together with the logarithms of the coefficients.
\begin{center}\rule{.18\linewidth}{.4pt}\end{center}
\par\vfill\noindent\hfill43\textsuperscript{*}
